\begin{theorem}[Young’s Inequality]
Let \(a,b\ge 0\) and \(p,q>1\) satisfy  
\[
\frac{1}{p}+\frac{1}{q}=1 .
\]
Then  
\[
ab\le \frac{a^{p}}{p}+\frac{b^{q}}{q}.
\]
Equality holds iff either \(a=0\), \(b=0\), or \(a>0,\;b>0\) and  
\[
a^{p}=b^{q}\qquad\bigl(\text{equivalently } a^{p-1}=b^{q-1}\bigr).
\]
\end{theorem}

\begin{proof}
\textbf{1. Conjugate exponents.}  
From \(\frac{1}{p}+\frac{1}{q}=1\) we obtain  
\[
q=\frac{p}{p-1},\qquad p=\frac{q}{q-1},
\]
so that indeed \(p,q\in(1,\infty)\).

\medskip

\textbf{2. Degenerate cases.}  
If \(a=0\) then \(ab=0\) and the right–hand side equals \(\frac{b^{q}}{q}\ge0\); thus the inequality holds, with equality only when also \(b=0\).  
The same argument works when \(b=0\).  
Henceforth assume \(a>0\) and \(b>0\).

\medskip

\textbf{3. Weighted AM–GM.}  
Set  
\[
\lambda_{1}:=\frac{1}{p},\qquad \lambda_{2}:=\frac{1}{q},
\]
so that \(\lambda_{1},\lambda_{2}>0\) and \(\lambda_{1}+\lambda_{2}=1\).  
The weighted arithmetic–geometric mean inequality states that for any non‑negative numbers \(x_{1},x_{2}\),
\begin{align}
\lambda_{1}x_{1}+\lambda_{2}x_{2}\ge x_{1}^{\lambda_{1}}x_{2}^{\lambda_{2}},
\tag{★}
\end{align}
with equality iff \(x_{1}=x_{2}\).  
Choosing \(x_{1}=a^{p}\) and \(x_{2}=b^{q}\) gives
\begin{align*}
\frac{1}{p}a^{p}+\frac{1}{q}b^{q}
&\ge (a^{p})^{1/p}\,(b^{q})^{1/q} \\
&= a\,b .
\end{align*}
Rearranging yields Young’s inequality
\[
ab\le\frac{a^{p}}{p}+\frac{b^{q}}{q}.
\]

\medskip

\textbf{4. Equality condition.}  
In \eqref{★} equality holds precisely when \(x_{1}=x_{2}\), i.e. when  
\[
a^{p}=b^{q},
\]
which is exactly the condition stated in the theorem, together with the degenerate cases treated in step 2.

\medskip

\textbf{5. Conclusion.}  
All hypotheses have been used; the desired inequality follows from the weighted AM–GM inequality, and the equality cases follow from the equality condition of that inequality. \(\qed\)
\end{proof}